\documentclass[french,11pt]{book}
\usepackage{teach}
\usepackage[french]{babel}
\usepackage{amsmath,amssymb}
\usepackage{array}
\newcommand{\R}{\mathbb{R}}
\begin{document}
\begin{center}
{\Large \textbf{Exercices -- Variables aléatoires}}\\[2mm]
Terminale technologique
\end{center}

\section*{Exercices}
\begin{enumerate}
  \item Une variable aléatoire $X$ prend les valeurs 0, 5 et 10 avec les probabilités 0,2 ; 0,5 et 0,3. Vérifier que la loi est complète puis calculer $E(X)$.
  \item Un jeu coûte 2 euros. On gagne 8 euros avec une probabilité de 0,15, 3 euros avec une probabilité de 0,25, et rien sinon. Définir le gain algébrique et calculer son espérance.
  \item Compléter une loi de probabilité sachant que $P(X=1)=0,18$, $P(X=2)=0,42$ et $P(X=3)=p$.
  \item Calculer la variance et l'écart type de la variable $X$ prenant les valeurs 1, 2, 4 avec probabilités 0,3 ; 0,5 ; 0,2.
\end{enumerate}

\end{document}
